\textbf{1. Kiến thức cần nhớ}
\begin{itemize}
\item \textbf{Góc nội tiếp} là góc có đỉnh nằm trên đường tròn và hai cạnh chứa hai dây cung của đường tròn đó.
\item \textbf{Định lí}: Trong một đường tròn, số đo của góc nội tiếp bằng \textbf{nửa số đo của cung bị chắn}:
$$\boxed{\widehat{BAC} = \dfrac{1}{2}\text{sđ}\wideparen{BC} = \dfrac{1}{2}\widehat{BOC}}$$
\item \textbf{Hệ quả quan trọng}: Trong một đường tròn:
\begin{itemize}
\item Các góc nội tiếp cùng chắn một cung (hoặc chắn các cung bằng nhau) thì \textbf{bằng nhau}.
\item Góc nội tiếp chắn \textbf{nửa đường tròn là góc vuông} ($\widehat{AMB} = 90^\circ$ khi $AB$ là đường kính).
\item Góc nội tiếp nhỏ hơn hoặc bằng $90^\circ$ có số đo bằng nửa góc ở tâm cùng chắn một cung.
\end{itemize}
\end{itemize}

\textbf{2. Phương pháp giải các dạng}

%% DANG 9H9?27-1
\textbf{Dạng 1. Tính số đo góc nội tiếp}
\begin{enumerate}
\item Xác định cung bị chắn bởi góc nội tiếp.
\item Áp dụng công thức $\text{Góc nội tiếp} = \dfrac{1}{2}\text{sđ(cung bị chắn)} = \dfrac{1}{2}\text{Góc ở tâm}$.
\item Vận dụng góc nội tiếp chắn nửa đường tròn luôn bằng $90^\circ$ để tìm các tam giác vuông.
\end{enumerate}

%% DANG 9H9?27-2
\textbf{Dạng 2. Chứng minh góc bằng nhau, hệ thức nhờ góc nội tiếp}
\begin{enumerate}
\item Chứng minh hai góc bằng nhau: chỉ ra chúng là hai góc nội tiếp cùng chắn một cung (hoặc hai cung bằng nhau).
\item Chứng minh tam giác đồng dạng theo trường hợp (g.g) nhờ các cặp góc nội tiếp bằng nhau, từ đó suy ra các hệ thức tích độ dài (như $MA \cdot MB = MC \cdot MD$).
\end{enumerate}
